\section{Conclusion}
\label{sec:conclusion}

We evaluated the multimodal in-context imitation learning capabilities of some of the world's most advanced AI models on interactive decision-making tasks --- tasks that are simple for humans but challenging for state-of-the-art LMs.
Our results show that, even with hundreds of demonstration episodes, context lengths of up to one million tokens, and thousands of output tokens, models often struggle to reach expert performance, thereby failing to translate their (factual) knowledge about the tasks' solution strategies into effective decision-making.
Solving this problem will be crucial for the next generational leap in LM capabilities towards general agents.
Despite our focus on in-context imitation learning, we believe it will be interesting to compare against other methods, including fine-tuning, retrieval-based methods, reward-conditioning, etc.
Our open-source benchmark (\url{https://github.com/google-deepmind/lm_act}) serves as a yard stick to measure progress toward that goal, and we are excited to see which innovations will be needed to solve all our simple tasks.

